% !TeX root = ../navier-stokes-manuscript.tex
% ============================================================
% 1.5 Class Membership and Participation
% ============================================================
\subsection{Class Membership and Participation}\label{sub:ClassMembershipAndParticipation}

The principle Silver needs has to answer one concrete question about every imagined failure: is this still the fluid generated by the original problem?
That question is class membership.

Fix the initial datum \(u_0\), the viscosity \(\nu>0\), and the maximal classical solution \((u,p)\) they generate on \([0,T_*)\).
Let \(Q=(\widetilde u,\widetilde p)\) be a proposed extension of this fixed solution to \([0,T_*+\varepsilon)\).
Write \(\operatorname{Member}(Q)\) when \(Q\) agrees with \((u,p)\) before \(T_*\) and continues to satisfy
\[
  \partial_t\widetilde u
  +(\widetilde u\cdot\nabla)\widetilde u
  +\nabla\widetilde p
  =
  \nu\Delta\widetilde u,
  \qquad
  \nabla\cdot\widetilde u=0,
\]
with the same datum, the same positive viscosity, the same whole-space decay and finite-energy requirements, and pressure still determined by the same velocity field.
When \(Q\) is not smooth enough for the classical derivatives, these requirements are read in the distributional and energy-inequality senses appropriate to the proposed class.
Smoothness is deliberately absent from membership, because it is the conclusion at issue.

Class exit is the complement
\[
  \operatorname{Exit}(Q)
  :=
  \neg\operatorname{Member}(Q).
\]
Now put the proposed extensions inside one ambient collection \(\mathcal Q\), and define
\[
  \mathcal M
  :=
  \{Q\in\mathcal Q:\operatorname{Member}(Q)\},
  \qquad
  \mathcal S
  :=
  \{Q\in\mathcal Q:\operatorname{Smooth}(Q)\}.
\]
The Gold Standard starts inside \(\mathcal M\) and proves that the admitted solution lies in \(\mathcal S\).
The Silver Standard starts outside \(\mathcal S\) and proves that the proposed failure lies outside \(\mathcal M\).
Both routes assert the same boundary:
\[
  \mathcal M\subseteq\mathcal S,
  \qquad
  \mathcal M\cap\mathcal S^{\mathsf c}=\varnothing.
\]
In implication form, this is
\[
  \operatorname{Member}(Q)
  \Longrightarrow
  \operatorname{Smooth}(Q),
  \qquad
  \neg\operatorname{Smooth}(Q)
  \Longrightarrow
  \operatorname{Exit}(Q).
\]
The Venn picture fixes the direction of the implication.
For Silver, each proposed nonsmooth outcome still needs a concrete reason it cannot remain a member.
A large derivative or a violent-looking picture counts only when it breaks one of the relations that makes the field the original Navier--Stokes fluid.

\divider{What Participation Means}

Those relations become physical when we follow the matter instead of holding the coordinate frame still.
Let \(X(a,t)\) be the trajectory of the parcel labelled by \(a\), so
\[
  \frac{d}{dt}X(a,t)
  =
  u(X(a,t),t).
\]
The parcel is moved by the velocity field and carried into a new neighbourhood of that same field.
Its velocity changes according to
\[
  \frac{d}{dt}u(X(a,t),t)
  =
  D_tu(X(a,t),t)
  =
  -\nabla p(X(a,t),t)
  +
  \nu\Delta u(X(a,t),t),
\]
where \(D_t=\partial_t+u\cdot\nabla\).

The two sides describe one event.
The parcel accelerates, while the pressure generated by the incompressibility of the whole field and the viscous comparison with its local velocity neighbourhood determine that acceleration.
Transport changes which spatial neighbourhood the parcel occupies and how nearby material is arranged, so the pressure and viscous responses are determined again from the changed field at the next instant.
No parcel supplies its own detached motion while the rest of the fluid evolves around it.

This continuing agreement is what participation means here.
A parcel participates while its position is transported by \(u\), its velocity responds through the pressure and viscous terms generated by \(u\), and its material neighbourhood remains inside the same divergence-free motion.
A small packet participates when that relation holds through the packet and the stresses across its moving boundary agree with the surrounding fluid.
Class membership says directly that this participation still belongs to the original problem: the same datum fixes the fluid, the same equations keep every parcel in the coupled motion, and the same pressure field keeps the local responses joined to the whole.

Participation does not require every reading to be nonzero.
A parcel can be instantaneously motionless while its acceleration is nonzero, and it can have zero acceleration when pressure and viscosity balance.
The zero solution has a vanishing derivative tower and remains a complete solution.
A steady flow has no Eulerian time change while its spatial structure continues to satisfy the momentum equation.
What matters is whether the value at each place agrees with the response demanded by the rest of the field.

To see that distinction instead of naming it abstractly, take one packet and prescribe two different tangential velocities across it.
Couette flow shows how one fluid can join those readings continuously, layer by neighbouring layer.
